% 原书第201—203页；10.3与10.3.1—10.3.2。
\section{学习因果时间序列模型}
目前，Granger因果性及其变体，是因果时间序列分析中最常用的方法之一。不过，为了更好地衔接各章，我们先介绍基于条件独立的方法能够得出哪些结论。这一介绍顺序绝不意味着对不同方法作出优劣评价。

第10.3.1节和第10.3.2节主要讨论可识别性结果。余下的第10.3.3—10.3.5节介绍更具体的时间序列因果学习方法。即使只在有限个时刻对多元时间序列进行过一次采样，也可以应用这些方法。不过，其中大多数思想也适用于对同一时间序列获得多次独立同分布重复观测的情况。

\subsection{马尔可夫条件与忠实性}
引理6.25指出，两个DAG马尔可夫等价，当且仅当它们具有相同骨架和相同的v结构集合。如果不存在即时效应，只要知道骨架，就已经确定了完整时间图，因为箭头只能沿时间向前。因此可以得出以下结论［Peters et al., 2013，定理1的证明］。
\begin{theorem}[不存在即时效应时的可识别性]
假设两个完整时间图分别由不存在即时效应的结构因果模型诱导。如果这两个完整时间图马尔可夫等价，则二者相同。
\end{theorem}

因此，只要马尔可夫条件与忠实性成立，就能从条件独立关系唯一识别完整时间图。为了处理包含无穷多个节点的DAG，有时会假设时间序列从$t=0$开始。

存在即时效应时，马尔可夫等价图至多在这些效应的方向上不同。不过，在许多情况下，即便这一方向也能识别，因为即时效应的不同方向会诱导不同的v结构集合。图10.5给出了一个简单例子。
\begin{figure}[htbp]
\centering
\begin{minipage}[t]{.47\linewidth}
\centering
\begin{tikzpicture}[causal,x=1.3cm,y=1.15cm]
\node[ch10node] (ch10f5a-X-0) at (0,1) {$X_{t-3}$};
\node[ch10node] (ch10f5a-X-1) at (1,1) {$X_{t-2}$};
\node[ch10node] (ch10f5a-X-2) at (2,1) {$X_{t-1}$};
\node[ch10node] (ch10f5a-X-3) at (3,1) {$X_{t}$};
\node[ch10node] (ch10f5a-Y-0) at (0,0) {$Y_{t-3}$};
\node[ch10node] (ch10f5a-Y-1) at (1,0) {$Y_{t-2}$};
\node[ch10node] (ch10f5a-Y-2) at (2,0) {$Y_{t-1}$};
\node[ch10node] (ch10f5a-Y-3) at (3,0) {$Y_{t}$};
\draw[ed] (ch10f5a-X-0)--(ch10f5a-X-1);
\draw[ed] (ch10f5a-X-1)--(ch10f5a-X-2);
\draw[ed] (ch10f5a-X-2)--(ch10f5a-X-3);
\draw[ed] (ch10f5a-Y-0)--(ch10f5a-Y-1);
\draw[ed] (ch10f5a-Y-1)--(ch10f5a-Y-2);
\draw[ed] (ch10f5a-Y-2)--(ch10f5a-Y-3);
\draw[ed] (ch10f5a-X-0)--(ch10f5a-Y-0);
\draw[ed] (ch10f5a-X-1)--(ch10f5a-Y-1);
\draw[ed] (ch10f5a-X-2)--(ch10f5a-Y-2);
\draw[ed] (ch10f5a-X-3)--(ch10f5a-Y-3);
\end{tikzpicture}
\par\small (a) $(Y_t)_{t\in\mathbb Z}$ 的所有节点处均存在 v 结构。
\end{minipage}\hfill
\begin{minipage}[t]{.47\linewidth}
\centering
\begin{tikzpicture}[causal,x=1.3cm,y=1.15cm]
\node[ch10node] (ch10f5b-X-0) at (0,1) {$X_{t-3}$};
\node[ch10node] (ch10f5b-X-1) at (1,1) {$X_{t-2}$};
\node[ch10node] (ch10f5b-X-2) at (2,1) {$X_{t-1}$};
\node[ch10node] (ch10f5b-X-3) at (3,1) {$X_{t}$};
\node[ch10node] (ch10f5b-Y-0) at (0,0) {$Y_{t-3}$};
\node[ch10node] (ch10f5b-Y-1) at (1,0) {$Y_{t-2}$};
\node[ch10node] (ch10f5b-Y-2) at (2,0) {$Y_{t-1}$};
\node[ch10node] (ch10f5b-Y-3) at (3,0) {$Y_{t}$};
\draw[ed] (ch10f5b-X-0)--(ch10f5b-X-1);
\draw[ed] (ch10f5b-X-1)--(ch10f5b-X-2);
\draw[ed] (ch10f5b-X-2)--(ch10f5b-X-3);
\draw[ed] (ch10f5b-Y-0)--(ch10f5b-Y-1);
\draw[ed] (ch10f5b-Y-1)--(ch10f5b-Y-2);
\draw[ed] (ch10f5b-Y-2)--(ch10f5b-Y-3);
\draw[ed] (ch10f5b-Y-0)--(ch10f5b-X-0);
\draw[ed] (ch10f5b-Y-1)--(ch10f5b-X-1);
\draw[ed] (ch10f5b-Y-2)--(ch10f5b-X-2);
\draw[ed] (ch10f5b-Y-3)--(ch10f5b-X-3);
\end{tikzpicture}
\par\small (b) $(X_t)_{t\in\mathbb Z}$ 的所有节点处均存在 v 结构。
\end{minipage}
\caption{两个DAG仅在即时效应上不同，却不马尔可夫等价。}
\end{figure}


即使在图10.5中对所有$t\in\mathbb Z$加入$X_t\to Y_{t+1}$，仍可推断即时效应的方向；改为加入$Y_t\to X_{t+1}$也一样。但不能同时加入这两类箭头，否则所有v结构都会消失。Peters et al.［2013，定理1］给出了即时效应方向可识别的下列充分条件。
\begin{theorem}[汇总图无环时的可识别性]
假设两个完整时间图由结构因果模型诱导，并且在两种情况下，对每个$j$，均存在某个$s\ge1$，使$X_t^j$受到$X_{t-s}^j$的影响。进一步假设汇总图无环。如果这两个完整时间图马尔可夫等价，则二者相同。
\end{theorem}

下面的结果说明，原则上只需一次条件独立性检验，就能判断汇总图中某条箭头是否存在。
\begin{theorem}[Granger因果性的依据]
考虑时间序列$(\mathbf X_t)_{t\in\mathbb Z}$上一个不存在即时效应的结构因果模型，并假设其诱导的联合分布相对于对应完整时间图满足忠实性。则汇总图包含从$X^j$到$X^k$的箭头，当且仅当存在$t\in\mathbb Z$使
\begin{equation}
 X_t^k\nindep X_{\mathrm{past}(t)}^j\mid\mathbf X_{\mathrm{past}(t)}^{-j}.
 \tag{10.3}
\end{equation}
\end{theorem}
为求完整，我们将证明收入附录C.14。类似结果可见White and Lu［2010］及Eichler［2011, 2012］。正如定理10.3的标题所提示的，这正是Granger因果性的基础，第10.3.3节将作更详细介绍。

\subsection{某些因果结论不需要忠实性}
值得关注的是，即使不用忠实性，也能从条件相依关系中得出有意义的因果结论。这不同于独立同分布情形：在后者中，无论如何排列节点，任意分布都相对于相应的完全DAG满足马尔可夫性。由于不存在沿时间向后的箭头，只要已知汇总图必为$X\to Y$或$Y\to X$之一，时间序列的马尔可夫条件就足以判断究竟是哪一种。

\Needspace{4cm}
\begin{theorem}[检测箭头$X\to Y$]
考虑二元时间序列$(X_t,Y_t)_{t\in\mathbb Z}$的一个结构因果模型。
\begin{enumerate}[label=(\roman*)]
 \item 如果存在某个$t\in\mathbb Z$使
 \begin{equation}
 Y_t\nindep X_{\mathrm{past}(t)}\mid Y_{\mathrm{past}(t)},
 \tag{10.4}
 \end{equation}
 那么汇总图包含从$X$到$Y$的箭头。
 \item 进一步假设不存在即时效应，且任意有限变量子集的联合密度均严格为正。如果对所有$t\in\mathbb Z$都有
 \begin{equation}
 Y_t\indep X_{\mathrm{past}(t)}\mid Y_{\mathrm{past}(t)},
 \tag{10.5}
 \end{equation}
 那么汇总图不包含从$X$到$Y$的箭头。
\end{enumerate}
\end{theorem}

这一证明可能也已在其他文献中出现；为求完整，我们将其收入附录C.15。证明第(ii)项需要因果极小性，而它严格弱于忠实性。
\chtennote{原书定理10.4的第(ii)项没有显式列出因果极小性假设，但紧接其后的文字及附录C.15的证明都要求这一条件。这里保留原陈述，应用该项结论时应同时考虑因果极小性。}

下一小节将说明，定理10.4及其多种变体［例如White and Lu, 2010；Eichler, 2011, 2012］，如何将基于条件独立的因果发现方法与Granger因果性联系起来。
\FloatBarrier
